\section{Background and Problem Setting}

\subsection{Entity Matching}

Entity matching is the task of deciding whether two descriptions from different data sources refer to the same real-world entity. It is a central step in data integration pipelines, because downstream tasks such as price comparison, product catalog construction, knowledge graph construction, and bibliographic data integration depend on consistent entity identifiers \cite{christophides2020overview,christophides2015webdata}. In the pairwise formulation used throughout this paper, each candidate pair $(l, r)$ consists of one record from a left table and one record from a right table. A label $y \in \{0,1\}$ indicates whether the two records match.

Most practical EM pipelines separate candidate construction from pair classification. Candidate construction, often called blocking, reduces the Cartesian product of two tables to a smaller set of plausible pairs. Pair classification then uses rules, similarity features, supervised classifiers, or neural models to decide whether each candidate pair is a match. Classical approaches include probabilistic record linkage \cite{fellegi1969record}, rule-based duplicate detection, and supervised feature-based systems \cite{elmagarmid2007duplicate,christen2012data}. Workflow-oriented systems such as Magellan highlight that EM is not only a classifier, but an end-to-end engineering process involving data preparation, blocking, labeling, debugging, and evaluation \cite{konda2016magellan}.

Neural EM methods replaced hand-crafted similarity features with learned representations. Early deep matching systems explored distributed tuple representations and neural pair classifiers \cite{mudgal2018deep}. Later systems adopted transformer encoders and pretrained language models (PLMs), such as BERT or RoBERTa, to serialize records into text and classify record pairs \cite{brunner2020transformer,li2020ditto}. Ditto is a representative PLM-based EM system that combines pretrained encoders with EM-specific data augmentation and domain knowledge injection \cite{li2020ditto}. These systems reach strong benchmark performance, but they still require task-specific labeled pairs and can be brittle when test pairs contain entities that were not observed during training \cite{peeters2025llmem}.

\subsection{LLM-Assisted Auto-Labeling}

Large language models (LLMs) provide a different interface to EM. Instead of fine-tuning a compact encoder on thousands of labeled pairs, a model can be prompted with two serialized entity descriptions and asked whether they refer to the same entity. Prior work by Peeters, Steiner, and Bizer showed that LLMs can perform competitively in zero-shot and few-shot EM settings, but also that there is no universally best prompt: prompt design and model choice interact strongly with the dataset \cite{peeters2025llmem}. The same work showed that LLMs can produce structured explanations for matching decisions and use these explanations to support error analysis.

Follow-up work on fine-tuning LLMs for EM showed that fine-tuning improves smaller models and in-domain transfer, while cross-domain generalization remains mixed \cite{steiner2025finetuning}. That study also explored training-example representation, including textual and structured explanations, and LLM-assisted example filtering and selection. These results motivate a closer look at the selected-and-labeled training data itself: if LLMs are used to label selected examples, the resulting labels should be treated as data artifacts with measurable provenance and quality.

The auto-labeling setting considered in this paper starts from a source benchmark and constructs a new labeled training set from pairs drawn from the two source tables. A pipeline first builds candidate pairs using retrieval or similarity search, then selects the most useful pairs for labeling, and finally asks an LLM to classify each selected pair. The output is a labeled dataset that can be used to train or evaluate downstream matchers.

This setting differs from direct LLM inference on a fixed test set. Direct inference asks whether an LLM can solve EM cases one at a time. Auto-labeling asks whether selected and LLM-labeled pairs form a training artifact with enough quality, coverage, and provenance for reliable downstream use. The latter requires auditing the selected-and-labeled data itself.

\subsection{Benchmark Setting}

The experiments in this draft follow the benchmark families used in the prior LLM-for-EM studies: product matching datasets such as WDC Products, Abt-Buy, Walmart-Amazon, and Amazon-Google, and bibliographic datasets such as DBLP-ACM and DBLP-Scholar \cite{peeters2025llmem,steiner2025finetuning}. These datasets expose different failure modes. Product matching often contains sparse or noisy descriptions, variants of the same product, and hard non-matches with similar surface forms. Bibliographic matching contains author, title, venue, and year variants, where small textual differences can be either harmless formatting variation or evidence for distinct publications.
